% !TeX root = ../main.tex
\chapter{State of the Art in Whole-Body Motion Planning for Mobile
  Manipulators}
\label{ch:sota}

The previous chapter argued that the Line-less Mobile Assembly System
(LMAS) setting calls for whole-body motion planning rather than a
decoupled base--arm pipeline. This chapter surveys how the literature
approaches the problem and narrows the field down to the framework used
in the remainder of this thesis. \Cref{sec:sota-classification}
introduces a classification of motion-planning approaches for mobile
manipulators. \Cref{sec:sota-holistic} then focuses on \emph{holistic}
(whole-body) planners, presenting the \ocs{} framework
(\cref{sec:sota-ocs2}) and the main alternatives
(\cref{sec:sota-alternatives}). \Cref{sec:sota-selection} compares the
candidates against the requirements of the \hero{} platform and
justifies the selection of \ocs{}. The synthesis of the applied and
open-source landscape draws, in part, on a systematic literature review
of whole-body motion planning and control integration in robotic
navigation stacks carried out within the same research
project~\cite{SantAnna2026WBMPReview}.

% ---------------------------------------------------------------------
\section{Classification of Motion Planning Approaches}
\label{sec:sota-classification}

Motion planning for mobile manipulators can be organised along a few
largely orthogonal axes. Systematic reviews of the
field~\cite{Sandakalum2022Survey,Thakar2023MobileManip} distinguish
approaches primarily by \emph{how the base and the arm are coupled} and
by \emph{the algorithmic principle} used to generate motion; a third,
practical axis concerns \emph{where the computation happens in time}.

\paragraph{Decoupled versus whole-body.}
The first and, for this work, most important distinction is between
\emph{decoupled} and \emph{whole-body} schemes.
A decoupled scheme plans the base and the arm separately and in
sequence: a navigation stack drives the base to a stand-off pose, after
which an arm planner solves the manipulation task from that fixed base
configuration. A whole-body scheme instead treats the base and the arm
as a \emph{single} redundant kinematic (or dynamic) system and solves for
their motion jointly, so that the task-space objective is distributed
between base and arm by the planner itself. As motivated in
\cref{sec:need-wholebody}, decoupling sacrifices reach, optimality, and
reactivity, which is precisely what whole-body planning recovers---at the
cost of a higher-dimensional and more tightly coupled problem.

\paragraph{Algorithmic principle.}
Independently of the coupling, the motion itself can be produced by
different families of algorithms:
\begin{itemize}
  \item \textbf{Sampling-based planning} explores the configuration space
        by randomly sampling and connecting collision-free states (for
        example Rapidly-exploring Random Trees and Probabilistic
        Roadmaps, as exposed through MoveIt~\cite{moveit2}). These methods
        excel at finding feasible paths in cluttered, high-dimensional
        spaces but do not natively optimise a running cost and are
        typically used offline.
  \item \textbf{Optimization-based planning} formulates motion as a
        constrained optimisation over a trajectory. When the optimisation
        is solved once, offline, over a full trajectory, one speaks of
        \emph{trajectory optimization}; when it is re-solved online over a
        receding horizon, the same principle becomes \emph{Model
        Predictive Control} (\mpc{})~\cite{Rawlings2017MPC}. Both handle
        joint, velocity, and collision constraints explicitly and produce
        smooth, cost-optimal motion.
  \item \textbf{Reactive schemes} compute an instantaneous command from
        the current state, for example through potential fields or
        quadratic-program-based inverse kinematics (QP-IK). They are cheap
        and fast but myopic, lacking the look-ahead needed to respect
        long-horizon constraints.
  \item \textbf{Learning-based schemes} replace or augment the planner
        with a policy learned by (deep) reinforcement learning. They can
        amortise the online cost of whole-body planning but raise
        questions of data efficiency, safety, and sim-to-real
        transfer~\cite{SantAnna2026WBMPReview}.
\end{itemize}

\paragraph{Offline versus online.}
Finally, a planner may operate \emph{offline}---computing a complete
trajectory before execution---or \emph{online}, re-planning continuously
as new state estimates and targets arrive. Only online, feedback-driven
planning provides the disturbance rejection required in a shared,
changing LMAS workspace. Whole-body \mpc{} sits at the intersection of the
\emph{whole-body}, \emph{optimization-based}, and \emph{online}
categories, which is the region of the design space this thesis targets.

% ---------------------------------------------------------------------
\section{Overview of Holistic Motion Planners}
\label{sec:sota-holistic}

\emph{Holistic} (whole-body) planners solve for base and arm motion in a
single formulation. In practice, the mature, deployable representatives of
this class are almost exclusively optimization-based, because expressing
the task in a cost function and the physical restrictions as constraints
maps naturally onto the coupled, redundant, and constrained nature of the
mobile-manipulation problem. A systematic review of how such controllers
are integrated with robotic navigation stacks found that, among the
open-source projects that couple a navigation stack with a genuine
whole-body controller, optimal-control formulations---and the \ocs{}
toolbox in particular---recur as a common backbone, alongside
operational-space control, quadratic-program solvers, and, more recently,
learned policies~\cite{SantAnna2026WBMPReview}. The following subsections
describe \ocs{} and then position the principal alternatives against it.

\subsection{The \ocs{} Framework: Optimal-Control-Based Model Predictive
  Control}
\label{sec:sota-ocs2}

\ocs{} (\emph{Optimal Control for Switched Systems}) is an open-source
C++ library for numerical optimal control developed at the Robotic
Systems Lab, ETH Z\"urich~\cite{ocs2Toolbox}. It provides efficient
solvers for continuous-time optimal-control problems---most notably
Sequential Linear-Quadratic (SLQ) control, a constrained variant of
differential dynamic programming (DDP)---together with the supporting
infrastructure required to run them as an \mpc{} controller.

Three features make \ocs{} well suited to whole-body mobile
manipulation. First, it derives the cost and constraint derivatives that
the solver needs by \emph{automatic differentiation}: the user supplies
the model (kinematics, cost terms, constraints) symbolically, and \ocs{}
generates and compiles efficient derivative code via CppAD/CppADCodeGen.
Second, it exposes a \emph{reference manager} that lets the desired
task-space trajectory be updated online, which is exactly the interface a
target-generating node needs in order to steer the controller. Third, and
central to a real-time deployment, \ocs{} separates the optimisation from
the execution through an MPC--MRT split: the \mpc{} process solves the
optimal-control problem at its own rate and publishes a \emph{policy},
while a lightweight Model Reference Tracking (MRT) client evaluates that
policy at the control rate to produce the instantaneous command. This
split is what makes it possible to run the optimiser and the hardware
interface as separate processes---even on separate machines---a property
this thesis exploits in the distributed architecture of
\cref{ch:integration}.

\ocs{} has been demonstrated for whole-body loco-manipulation on legged
and wheeled platforms, including a unified \mpc{} formulation for combined
locomotion and manipulation~\cite{Sleiman2021UnifiedMPC} and perceptive
whole-body interaction with articulated
objects~\cite{Mittal2022Articulated}. Its lineage also includes
foundational work on kinematic planning for non-holonomic mobile
manipulators and whole-body \mpc{} for dynamically stable mobile
manipulators, which the review in the same project identifies as the core
mobile-manipulation references behind the
toolbox~\cite{SantAnna2026WBMPReview}. This track record on hardware, its
explicit constraint handling, and its process-separable architecture make
it the primary candidate for the \hero{} platform.

\subsection{Alternative Frameworks and Methods}
\label{sec:sota-alternatives}

Several alternatives to \ocs{} exist, spanning the algorithmic families of
\cref{sec:sota-classification}.

\paragraph{Sampling-based planning (MoveIt).}
MoveIt~\cite{moveit2} is the de-facto standard motion-planning framework
in the ROS ecosystem. Its strength is offline, collision-free planning
for redundant arms in cluttered scenes; it can be extended to a mobile
base by treating the base pose as additional planning variables. In its
common usage, however, MoveIt underpins the \emph{decoupled}
navigate-then-manipulate pipeline and does not, by itself, provide the
online, coordinated base--arm feedback that whole-body \mpc{} offers. It
is therefore the natural baseline against which a whole-body approach is
compared rather than a competing whole-body planner.

\paragraph{Other optimization-based / \mpc{} tools.}
Beyond \ocs{}, whole-body behaviour can be realised with general-purpose
optimal-control and \mpc{} solvers or with constraint-based task-space
controllers. The project-internal review catalogues open-source stacks
that integrate such controllers with a navigation stack through different
mechanisms---replacing the Nav2 local planner via the
\texttt{nav2\_core::Controller} plugin interface, wrapping the controller
in \texttt{ros2\_control}, or bridging \texttt{cmd\_vel}/goal topics to a
whole-body solver---and reports that \ocs{}-based and quadratic-program
solvers dominate this category~\cite{SantAnna2026WBMPReview}. These tools
are technically capable but generally require more integration effort or
offer less mature support for the specific combination of a
differential-drive base and a ROS~2 hardware stack targeted here.

\paragraph{Reactive quadratic-program inverse kinematics.}
QP-based inverse kinematics resolves the \ee{} task and secondary
objectives (joint-limit avoidance, manipulability) into joint velocities
at each control step. It is computationally light and reactive but, being
a one-step method, cannot anticipate constraints over a horizon, which
limits its ability to plan coordinated base--arm motion for long,
reach-exceeding paths.

\paragraph{Learning-based whole-body control.}
A growing body of work hybridises classical navigation with reinforcement
learning to amortise the online cost of whole-body planning, learning
policies that map observations directly to base and arm
commands~\cite{SantAnna2026WBMPReview}. While promising, these approaches
introduce training-data, reproducibility, and sim-to-real-transfer
concerns that are not aligned with the goal of this thesis---deploying a
transparent, analysable controller on real hardware in the near term.

% ---------------------------------------------------------------------
\section{Comparison of Motion Planning Algorithms and Selection of a
  Suitable Algorithm}
\label{sec:sota-selection}

The choice of a planner for \hero{} is driven by the requirements
identified in \cref{ch:introduction}: coordinated whole-body motion,
explicit handling of joint/velocity, non-holonomic, and self-collision
constraints, online reactivity, real-time feasibility on modest on-board
compute, and good ROS~2 support so that the controller can be integrated
with the existing hardware stack. \Cref{tab:sota-comparison} summarises
how the candidate approaches score against these criteria.

\begin{table}[htbp]
  \centering
  \caption[Comparison of motion-planning approaches for \hero{}]{%
    Qualitative comparison of motion-planning approaches against the
    requirements of the \hero{} platform
    ($+$ good, $\circ$ partial, $-$ weak).}
  \label{tab:sota-comparison}
  \begin{tabularx}{\linewidth}{@{}Xcccc@{}}
    \toprule
    Criterion
      & \shortstack{Sampling\\(MoveIt)}
      & \shortstack{QP-IK\\(reactive)}
      & \shortstack{Learning\\(RL)}
      & \shortstack{Whole-body\\\mpc{} (\ocs{})} \\
    \midrule
    Coordinated whole-body motion & $\circ$ & $\circ$ & $+$      & $+$      \\
    Explicit constraint handling  & $\circ$ & $\circ$ & $-$      & $+$      \\
    Online reactivity             & $-$     & $+$     & $+$      & $+$      \\
    Real-time feasibility         & $\circ$ & $+$     & $+$      & $\circ$  \\
    Redundancy resolution         & $\circ$ & $+$     & $+$      & $+$      \\
    ROS~2 maturity                & $+$     & $\circ$ & $-$      & $+$      \\
    Transparency / analysability  & $+$     & $+$     & $-$      & $+$      \\
    \bottomrule
  \end{tabularx}
\end{table}

Sampling-based planning is mature and well supported but is offline and,
in its standard form, decoupled---disqualifying it as the primary
whole-body planner while keeping it relevant as a baseline. Reactive
QP-IK is fast and reactive but myopic. Learning-based control is
attractive in principle but conflicts with the near-term, hardware-first,
and analysability goals of this work. Whole-body \mpc{} is the only family
that simultaneously delivers coordinated base--arm motion, explicit
constraint handling, and online reactivity; its main drawback---the
computational burden of solving an optimal-control problem online---is
addressable through the MPC--MRT split and a distributed deployment
(\cref{ch:integration}) and is quantified experimentally in
\cref{ch:hardware}.

Among whole-body \mpc{} implementations, \ocs{} is selected for its mature
open-source implementation, its explicit constraint support, its
demonstrated use on wheeled and legged mobile manipulators, and---decisive
for a real-time deployment on distributed hardware---its process-separable
MPC--MRT architecture. The remainder of this thesis therefore integrates
and validates \ocs{}-based whole-body \mpc{} on the \hero{} platform.
